\chapter{§2.9 不等式}
\label{sec:2.9}


\prerequisite{§2.6 一元一次方程}

\gradelevel{7-8 年级}


\section*{不等式的概念与基本性质}


\begin{explore}


超市购物时，小明带了 $50$ 元，想买若干瓶 $3$ 元的饮料。他不能花超过 $50$ 元，即 $3x \leq 50$ 。这不是一个"等于"的关系，而是一个"不超过"的关系——这就是\textbf{不等式}。生活中，很多问题不是精确相等的，而是涉及大小比较。

\end{explore}

\begin{understand}


\textbf{不等式}：用不等号（ $>$ ， $<$ ， $\geq$ ， $\leq$ ， $\neq$ ）连接的式子叫做不等式。

常见不等号的含义：

\begin{center}
\begin{tabular}{ll}
\toprule
符号 & 含义 \\
\midrule
$a > b$ & $a$ 大于 $b$ \\
$a < b$ & $a$ 小于 $b$ \\
$a \geq b$ & $a$ 大于或等于 $b$ \\
$a \leq b$ & $a$ 小于或等于 $b$ \\
\bottomrule
\end{tabular}
\end{center}


\textbf{不等式的基本性质}：

\textbf{性质 1}：不等式两边同时加上（或减去）同一个数或式子，不等号方向不变。

\[
\text{若 } a > b, \text{ 则 } a + c > b + c
\]


\textbf{性质 2}：不等式两边同时乘以（或除以）同一个\textbf{正数}，不等号方向不变。

\[
\text{若 } a > b, \; c > 0, \text{ 则 } ac > bc
\]


\textbf{性质 3}：不等式两边同时乘以（或除以）同一个\textbf{负数}，不等号方向\textbf{改变}。

\[
\text{若 } a > b, \; c < 0, \text{ 则 } ac < bc
\]


\textbf{为什么乘以负数要变号？} 直觉上： $3 > 2$ ，但 $3 \times (-1) = -3 < -2 = 2 \times (-1)$ 。乘以负数相当于在数轴上"翻转"，大小关系也随之翻转。

\end{understand}

\begin{workedexamples}


\textbf{例 1.} 用不等式表示： $x$ 是正数； $y$ 不超过 $5$ ； $a$ 与 $b$ 的和不小于 $10$ 。

\textbf{解：}

\begin{itemize}
  \item $x > 0$
  \item $y \leq 5$
  \item $a + b \geq 10$
\end{itemize}


\textbf{例 2.} 已知 $a > b$ ，判断下列不等式是否成立：

(1) $a + 3 > b + 3$ 　(2) $-2a > -2b$ 　(3) $\frac{a}{-5} > \frac{b}{-5}$

\textbf{解：}

(1) 两边加 $3$ ，不等号不变。✓ 成立。

(2) 两边乘 $-2$ （负数），不等号改变。 $-2a < -2b$ 。✗ 不成立。

(3) 两边除以 $-5$ （负数），不等号改变。 $\frac{a}{-5} < \frac{b}{-5}$ 。✗ 不成立。

\textbf{例 3.} 若 $-3a < -3b$ ，比较 $a$ 与 $b$ 的大小。

\textbf{解：} $-3a < -3b$ ，两边除以 $-3$ （负数），不等号改变： $a > b$ 。

\end{workedexamples}

\begin{keytakeaway}


\begin{itemize}
  \item 不等式的性质 1 和性质 2 与等式性质类似。
  \item 性质 3 是不等式独有的：乘除负数要变号，这是最关键也最易错的。
  \item 不等式性质是解不等式的理论基础。
\end{itemize}


\end{keytakeaway}

\begin{practice}


\begin{enumerate}
  \item 用不等式表示： $x$ 是非负数； $m$ 大于 $-3$ 。
  \item 已知 $a < b$ ，判断 $5a$ 与 $5b$ ， $-a$ 与 $-b$ 的大小关系。
  \item 若 $a > b > 0$ ，则 $\frac{1}{a}$ 与 $\frac{1}{b}$ 哪个大？
  \item 已知 $c < 0$ ， $a > b$ ，比较 $ac$ 与 $bc$ 。
  \item 不等式两边乘以 $0$ 会怎样？这一操作有什么问题？
\end{enumerate}


\end{practice}

\section*{一元一次不等式的解法}


\begin{explore}


解方程 $3x - 5 = 7$ 得到一个确定的值 $x = 4$ 。但如果改为 $3x - 5 > 7$ ，答案就不是一个值，而是一个范围：所有大于 $4$ 的数都满足条件。解不等式就是找出所有使不等式成立的未知数的值。

\end{explore}

\begin{understand}


\textbf{一元一次不等式}：含一个未知数，未知数的最高次数为 $1$ 的不等式。一般形式为 $ax + b > 0$ （或 $<, \geq, \leq$ ）。

\textbf{解不等式}的步骤与解方程完全类似：

\begin{enumerate}
  \item 去分母、去括号。
  \item 移项（变号规则同方程）。
  \item 合并同类项。
  \item 系数化为 $1$ ——\textbf{注意}：除以负数时，不等号要改变方向！
\end{enumerate}


\textbf{不等式的解}：使不等式成立的未知数的值。

\textbf{不等式的解集}：所有解组成的集合。

\end{understand}

\begin{workedexamples}


\textbf{例 1.} 解不等式 $2x - 3 > 5$ 。

\textbf{解：}

\[
2x > 5 + 3 = 8
\]


\[
x > 4
\]


解集为 $x > 4$ 。

\textbf{例 2.} 解不等式 $-3x + 6 \leq 0$ 。

\textbf{解：}

\[
-3x \leq -6
\]


两边除以 $-3$ （负数，变号）：

\[
x \geq 2
\]


解集为 $x \geq 2$ 。

\textbf{例 3.} 解不等式 $\frac{2x - 1}{3} > \frac{x + 2}{2}$ 。

\textbf{解：} 去分母（两边乘 $6$ ）：

\[
2(2x - 1) > 3(x + 2)
\]


\[
4x - 2 > 3x + 6
\]


\[
x > 8
\]


解集为 $x > 8$ 。

\textbf{例 4.} 解不等式 $3(x - 2) - 1 \leq 4(x + 1)$ 。

\textbf{解：}

\[
3x - 6 - 1 \leq 4x + 4
\]


\[
3x - 7 \leq 4x + 4
\]


\[
-x \leq 11
\]


\[
x \geq -11
\]


解集为 $x \geq -11$ 。

\end{workedexamples}

\begin{keytakeaway}


\begin{itemize}
  \item 解不等式的步骤与解方程基本相同。
  \item 唯一不同：两边乘除\textbf{负数}时，不等号方向要改变。
  \item 不等式的解是一个范围（解集），不是一个值。
\end{itemize}


\end{keytakeaway}

\begin{practice}


\begin{enumerate}
  \item 解不等式 $5x - 2 > 8$ 。
  \item 解不等式 $-4x + 3 < 11$ 。
  \item 解不等式 $2(x - 1) \geq 3x + 1$ 。
  \item 解不等式 $\frac{x + 1}{2} \leq \frac{3x - 1}{4}$ 。
  \item 解不等式 $1 - \frac{x - 3}{6} > \frac{2x}{3}$ 。
  \item 找出使不等式 $3x + 7 \leq 0$ 成立的最大整数。
\end{enumerate}


\end{practice}

\section*{不等式的解集与数轴表示}


\begin{explore}


不等式 $x > 3$ 的解集包含 $3.1, 4, 100, \ldots$ 无穷多个数，没法一一列出。但在数轴上，这些数恰好对应 $3$ 右边的所有点——一条射线！数轴表示法让我们可以直观地"看到"解集。

\end{explore}

\begin{understand}


不等式解集的\textbf{数轴表示}：

\begin{center}
\begin{tabular}{ll}
\toprule
解集 & 数轴表示 \\
\midrule
$x > a$ & $a$ 处\textbf{空心圆}（不含 $a$ ），向右画射线 \\
$x \geq a$ & $a$ 处\textbf{实心圆}（含 $a$ ），向右画射线 \\
$x < a$ & $a$ 处空心圆，向左画射线 \\
$x \leq a$ & $a$ 处实心圆，向左画射线 \\
\bottomrule
\end{tabular}
\end{center}


记忆口诀：\textbf{大于向右，小于向左；有等号用实心，无等号用空心}。

\end{understand}

\begin{workedexamples}


\textbf{例 1.} 在数轴上表示解集 $x > 2$ 。

\textbf{解：} 在数轴上 $2$ 的位置画空心圆，从 $2$ 向右画射线（含箭头）。

\[
\underset{2}{\circ} \longrightarrow
\]


\textbf{例 2.} 在数轴上表示解集 $x \leq -1$ 。

\textbf{解：} 在 $-1$ 的位置画实心圆，从 $-1$ 向左画射线。

\[
\longleftarrow \underset{-1}{\bullet}
\]


\textbf{例 3.} 解不等式 $4 - 2x > 0$ 并在数轴上表示解集。

\textbf{解：} $-2x > -4$ ， $x < 2$ 。

数轴上：在 $2$ 处画空心圆，向左画射线。

\textbf{例 4.} 解不等式 $3x - 1 \geq 2x + 3$ ，在数轴上表示，并写出满足条件的最小整数。

\textbf{解：} $x \geq 4$ 。数轴上 $4$ 处实心圆，向右射线。满足条件的最小整数是 $4$ 。

\end{workedexamples}

\begin{keytakeaway}


\begin{itemize}
  \item 数轴是不等式解集的可视化工具。
  \item 空心圆表示"不含"，实心圆表示"含"。
  \item 数轴表示对于后续学习不等式组非常重要。
\end{itemize}


\end{keytakeaway}

\begin{practice}


\begin{enumerate}
  \item 在数轴上表示 $x \geq -3$ 。
  \item 在数轴上表示 $x < 5$ 。
  \item 解不等式 $-x + 4 > 1$ ，在数轴上表示解集。
  \item 解不等式 $2x + 6 \leq 0$ ，在数轴上表示解集，并写出满足条件的最大整数。
  \item 写出 $-2 < x \leq 3$ 对应的数轴描述。
\end{enumerate}


\end{practice}

\section*{一元一次不等式组}


\begin{explore}


体育课上，老师要求跳远成绩至少 $2$ 米但不超过 $3$ 米的同学进入下一轮。设成绩为 $x$ 米，则要求同时满足 $x \geq 2$ 和 $x \leq 3$ ，即 $2 \leq x \leq 3$ 。两个不等式联合使用就是\textbf{不等式组}。

\end{explore}

\begin{understand}


\textbf{一元一次不等式组}：由两个（或多个）一元一次不等式组成的不等式组。

解不等式组的步骤：

\begin{enumerate}
  \item 分别求出每个不等式的解集。
  \item 在数轴上标出每个解集。
  \item 取\textbf{公共部分}（交集）。
\end{enumerate}


确定公共部分的口诀：

\begin{itemize}
  \item \textbf{同大取大}： $x > 3$ 且 $x > 5$ → $x > 5$
  \item \textbf{同小取小}： $x < 3$ 且 $x < 5$ → $x < 3$
  \item \textbf{大小小大取中间}： $x > 2$ 且 $x < 7$ → $2 < x < 7$
  \item \textbf{大大小小无解}： $x > 5$ 且 $x < 2$ → 无解
\end{itemize}


\end{understand}

\begin{workedexamples}


\textbf{例 1.} 解不等式组 $\begin{cases} x + 3 > 0 \\ 2x - 4 < 0 \end{cases}$ 。

\textbf{解：}

由 $x + 3 > 0$ 得 $x > -3$ 。

由 $2x - 4 < 0$ 得 $x < 2$ 。

公共部分： $-3 < x < 2$ 。

\textbf{例 2.} 解不等式组 $\begin{cases} 3x - 6 \geq 0 \\ 5x - 10 \leq 0 \end{cases}$ 。

\textbf{解：}

由 $3x - 6 \geq 0$ 得 $x \geq 2$ 。

由 $5x - 10 \leq 0$ 得 $x \leq 2$ 。

公共部分： $x = 2$ （唯一值）。

\textbf{例 3.} 解不等式组 $\begin{cases} 2x + 1 > 5 \\ 3x - 2 > 7 \end{cases}$ 。

\textbf{解：}

由 $2x + 1 > 5$ 得 $x > 2$ 。

由 $3x - 2 > 7$ 得 $x > 3$ 。

同大取大： $x > 3$ 。

\textbf{例 4.} 解不等式组 $\begin{cases} x - 1 > 3 \\ 2x + 5 < 1 \end{cases}$ 。

\textbf{解：}

由 $x - 1 > 3$ 得 $x > 4$ 。

由 $2x + 5 < 1$ 得 $x < -2$ 。

大大小小，公共部分为空集。不等式组\textbf{无解}。

\end{workedexamples}

\begin{keytakeaway}


\begin{itemize}
  \item 不等式组的解集是各不等式解集的交集。
  \item 用数轴找交集最直观。
  \item 口诀：同大取大、同小取小、大小小大取中间、大大小小无解。
\end{itemize}


\end{keytakeaway}

\begin{practice}


\begin{enumerate}
  \item 解不等式组 $\begin{cases} x - 2 > 0 \\ x + 1 < 6 \end{cases}$ 。
  \item 解不等式组 $\begin{cases} 3x \geq 6 \\ 2x \leq 10 \end{cases}$ 。
  \item 解不等式组 $\begin{cases} 2x - 1 > 3 \\ x + 4 > 8 \end{cases}$ 。
  \item 解不等式组 $\begin{cases} x + 5 > 0 \\ x - 3 > 0 \end{cases}$ ，并写出满足条件的最小正整数。
  \item 解不等式组 $\begin{cases} 3x - 1 \leq 8 \\ 2x + 3 \geq 1 \end{cases}$ ，并写出满足条件的所有整数。
  \item 解不等式组 $\begin{cases} 4x > 12 \\ -x > 2 \end{cases}$ 。
\end{enumerate}


\end{practice}

\section*{不等式的应用}


\begin{explore}


班级要印制活动手册，甲印刷厂报价：每本 $3$ 元，另收版费 $500$ 元；乙印刷厂报价：每本 $4.5$ 元，不收版费。印多少本以上时选甲厂更划算？

\end{explore}

\begin{understand}


不等式的应用与方程应用类似，但题目中通常出现"至少""至多""不超过""不低于"等表示范围的词语。列不等式时注意对应的不等号。

\begin{center}
\begin{tabular}{ll}
\toprule
关键词 & 对应的不等号 \\
\midrule
至少、不低于、不少于 & $\geq$ \\
至多、不超过、不多于 & $\leq$ \\
超过、大于 & $>$ \\
不足、小于 & $<$ \\
\bottomrule
\end{tabular}
\end{center}


\end{understand}

\begin{workedexamples}


\textbf{例 1.} 甲厂总费用 $3x + 500$ ，乙厂总费用 $4.5x$ 。甲比乙便宜的条件是：

\[
3x + 500 < 4.5x
\]


\[
500 < 1.5x
\]


\[
x > \frac{1000}{3} \approx 333.3
\]


所以印 $334$ 本以上时选甲厂更划算。

\textbf{例 2.} 某班组织郊游，租一辆大巴需 $800$ 元，每人分摊车费不超过 $20$ 元，至少需要多少人参加？

\textbf{解：} 设 $x$ 人参加。 $\frac{800}{x} \leq 20$ ， $800 \leq 20x$ ， $x \geq 40$ 。

至少需要 $40$ 人参加。

\textbf{例 3.} 小红期末考试前四次数学测验的成绩分别为 $85, 78, 92, 88$ 分。第五次至少考多少分才能使五次的平均分不低于 $85$ ？

\textbf{解：} 设第五次考 $x$ 分。

\[
\frac{85 + 78 + 92 + 88 + x}{5} \geq 85
\]


\[
343 + x \geq 425
\]


\[
x \geq 82
\]


至少考 $82$ 分。

\textbf{例 4.} 一根 $20$ 厘米的绳子要剪成两段，其中一段的长度不超过另一段的 $2$ 倍。求较短那段的取值范围。

\textbf{解：} 设较短段为 $x$ 厘米，较长段为 $(20 - x)$ 厘米。

条件： $x \leq 20 - x$ （ $x$ 为较短段）且 $20 - x \leq 2x$ （较长段不超过较短段的 $2$ 倍）。

由 $x \leq 20 - x$ ： $2x \leq 20$ ， $x \leq 10$ 。

由 $20 - x \leq 2x$ ： $20 \leq 3x$ ， $x \geq \frac{20}{3}$ 。

又因为 $x > 0$ ，所以 $\frac{20}{3} \leq x \leq 10$ 。

\end{workedexamples}

\begin{keytakeaway}


\begin{itemize}
  \item 应用题中注意关键词对应的不等号。
  \item 不等式的解通常是一个范围，需要根据实际情况取整数或合理值。
  \item 有些问题需要列不等式组（多个约束条件同时满足）。
\end{itemize}


\end{keytakeaway}

\begin{practice}


\begin{enumerate}
  \item 一个数的 $3$ 倍减去 $5$ 不小于 $10$ ，求这个数的取值范围。
  \item 某商店购进一批单价 $50$ 元的商品，售价至少定为多少才能使利润率不低于 $20\%$ ？
  \item 小明有 $100$ 元，买了若干本 $12$ 元的笔记本后还想留至少 $20$ 元。最多能买多少本？
  \item 一列火车从甲站到乙站，按时刻表应该在 $3$ 小时内到达（含 $3$ 小时），已知两站间距 $300$ 千米，火车平均速度至少是多少？
  \item 某次考试满分 $100$ 分，小王前两次分别考了 $72$ 分和 $85$ 分。要使三次平均分超过 $80$ 分，第三次至少考多少分？
\end{enumerate}


\end{practice}



\begin{answer}


\textbf{不等式的概念与基本性质 练一练}

\begin{enumerate}
  \item $x \geq 0$ ； $m > -3$ 。
  \item $a < b$ 两边乘 $5$ （正数）， $5a < 5b$ 。两边乘 $-1$ （负数）， $-a > -b$ 。
  \item 由 $a > b > 0$ ，两边取倒数，不等号改变： $\frac{1}{a} < \frac{1}{b}$ 。
  \item $c < 0$ ， $a > b$ 两边乘 $c$ ，变号： $ac < bc$ 。
  \item 两边乘 $0$ 得 $0 > 0$ （或 $0 < 0$ ），这是假命题。所以不等式两边不能乘以 $0$ 。
\end{enumerate}


\textbf{一元一次不等式的解法 练一练}

\begin{enumerate}
  \item $5x > 10$ ， $x > 2$ 。
  \item $-4x < 8$ ， $x > -2$ 。
  \item $2x - 2 \geq 3x + 1$ ， $-x \geq 3$ ， $x \leq -3$ 。
  \item 去分母（乘 $4$ ）： $2(x + 1) \leq 3x - 1$ ， $2x + 2 \leq 3x - 1$ ， $-x \leq -3$ ， $x \geq 3$ 。
  \item 去分母（乘 $6$ ）： $6 - (x - 3) > 4x$ ， $6 - x + 3 > 4x$ ， $9 > 5x$ ， $x < \frac{9}{5}$ 。
  \item $3x \leq -7$ ， $x \leq -\frac{7}{3} \approx -2.33$ 。最大整数为 $-3$ 。
\end{enumerate}


\textbf{不等式的解集与数轴表示 练一练}

\begin{enumerate}
  \item $-3$ 处实心圆，向右射线。
  \item $5$ 处空心圆，向左射线。
  \item $-x > -3$ ， $x < 3$ 。数轴上 $3$ 处空心圆，向左射线。
  \item $2x \leq -6$ ， $x \leq -3$ 。数轴上 $-3$ 处实心圆，向左射线。最大整数为 $-3$ 。
  \item $-2$ 处空心圆， $3$ 处实心圆，中间连线。表示大于 $-2$ 且小于或等于 $3$ 的所有数。
\end{enumerate}


\textbf{一元一次不等式组 练一练}

\begin{enumerate}
  \item $x > 2$ 且 $x < 5$ ，解集 $2 < x < 5$ 。
  \item $x \geq 2$ 且 $x \leq 5$ ，解集 $2 \leq x \leq 5$ 。
  \item $x > 2$ 且 $x > 4$ ，同大取大， $x > 4$ 。
  \item $x > -5$ 且 $x > 3$ ，同大取大， $x > 3$ 。最小正整数为 $4$ 。
  \item $3x \leq 9$ ， $x \leq 3$ ； $2x \geq -2$ ， $x \geq -1$ 。解集 $-1 \leq x \leq 3$ 。整数： $-1, 0, 1, 2, 3$ 。
  \item $x > 3$ 且 $x < -2$ ，大大小小，无解。
\end{enumerate}


\textbf{不等式的应用 练一练}

\begin{enumerate}
  \item $3x - 5 \geq 10$ ， $3x \geq 15$ ， $x \geq 5$ 。
  \item $\frac{x - 50}{50} \geq 20\%$ ， $x - 50 \geq 10$ ， $x \geq 60$ 。售价至少 $60$ 元。
  \item $12x \leq 100 - 20 = 80$ ， $x \leq 6.67$ 。最多买 $6$ 本。
  \item $\frac{300}{v} \leq 3$ ， $v \geq 100$ 。平均速度至少 $100$ 千米/时。
  \item $\frac{72 + 85 + x}{3} > 80$ ， $157 + x > 240$ ， $x > 83$ 。第三次至少考 $84$ 分（取最小整数）。
\end{enumerate}


\end{answer}
